%% 
%% Copyright 2007-2025 Elsevier Ltd
%% 
%% This file is part of the 'Elsarticle Bundle'.
%% ---------------------------------------------
%% 
%% It may be distributed under the conditions of the LaTeX Project Public
%% License, either version 1.3 of this license or (at your option) any
%% later version.  The latest version of this license is in
%%    http://www.latex-project.org/lppl.txt
%% and version 1.3 or later is part of all distributions of LaTeX
%% version 1999/12/01 or later.
%% 
%% The list of all files belonging to the 'Elsarticle Bundle' is
%% given in the file `manifest.txt'.
%% 
%% Template article for Elsevier's document class `elsarticle'
%% with harvard style bibliographic references

\documentclass[preprint,3p,12pt,authoryear]{elsarticle}

%% Use the option review to obtain double line spacing
%% \documentclass[authoryear,preprint,review,12pt]{elsarticle}

%% Use the options 1p,twocolumn; 3p; 3p,twocolumn; 5p; or 5p,twocolumn
%% for a journal layout:
%% \documentclass[final,1p,times,authoryear]{elsarticle}
%% \documentclass[final,1p,times,twocolumn,authoryear]{elsarticle}
%% \documentclass[final,3p,times,authoryear]{elsarticle}
%% \documentclass[final,3p,times,twocolumn,authoryear]{elsarticle}
%% \documentclass[final,5p,times,authoryear]{elsarticle}
%% \documentclass[final,5p,times,twocolumn,authoryear]{elsarticle}

%% For including figures, graphicx.sty has been loaded in
%% elsarticle.cls. If you prefer to use the old commands
%% please give \usepackage{epsfig}

%% The amssymb package provides various useful mathematical symbols
\usepackage{amssymb}
%% The amsmath package provides various useful equation environments.
\usepackage{amsmath}
\usepackage{float}
\usepackage{booktabs}
\usepackage{longtable}
%% The amsthm package provides extended theorem environments
%% \usepackage{amsthm}

%% The lineno packages adds line numbers. Start line numbering with
%% \begin{linenumbers}, end it with \end{linenumbers}. Or switch it on
%% for the whole article with \linenumbers.
%% \usepackage{lineno}

\journal{Expert Systems with Applications}

\begin{document}

\begin{frontmatter}

\title{A Three-Stage Optimization Framework for Remote-Sensing Satellite Mission Scheduling}

\begin{abstract}
Remote-sensing satellite mission scheduling requires coordinated decisions on image acquisition, data playback, and on-board memory release under tightly coupled temporal and resource constraints. Existing integrated scheduling studies usually model imaging and downloading, but do not explicitly represent the full capture-download-delete lifecycle that governs practical memory recovery. This paper presents a three-stage optimization framework for remote-sensing satellite mission scheduling that combines mathematical modeling with staged heuristic-exact decision support. We first formulate a mixed-integer model that captures lifecycle coupling, memory dynamics, ground-station coordination, and orbital duty-cycle limits. We then develop a three-stage solution architecture that combines mission feasibility selection, playback sequencing through a knapsack-to-shortest-path mechanism, and linear-programming-based timing refinement. The framework is designed to generate high-quality feasible schedules within operationally acceptable runtimes and to support practical planning decisions under realistic satellite-control constraints. Computational experiments on instances ranging from 24 to 167 missions show that the proposed framework consistently produces competitive schedules and outperforms time-limited CPLEX on medium and large instances in mission completion and computational efficiency. Sensitivity analyses further show how memory capacity, contact availability, and download speed affect scheduling performance and infrastructure requirements. The framework has been implemented in operational satellite-control settings, demonstrating its applicability as a practical decision-support tool for mission planning.
\end{abstract}

\begin{keyword}
Satellite mission scheduling \sep Three-stage optimization \sep Decision support \sep Hybrid heuristic \sep Mixed-integer programming \sep Remote sensing
\end{keyword}

\end{frontmatter}

%% Start line numbering for review
% \linenumbers

% INTRODUCTION
\section{Introduction}
Remote-sensing satellites have become indispensable assets, providing critical data for applications ranging from environmental monitoring to disaster management \citep{Madry2013}. These platforms orbit the Earth, collecting high-resolution surface information that is subsequently downlinked to ground stations for processing. However, as the demand for satellite data and the sophistication of satellite technology have grown, so too has the operational complexity of managing their missions. The core challenge lies in the efficient scheduling of diverse and often conflicting tasks under a set of stringent resource and temporal constraints. This has transformed satellite scheduling into a large-scale, NP-hard combinatorial optimization problem that demands increasingly advanced solutions.

This complexity is driven by technological advancement. The transition from early-generation platforms to modern Agile Earth Observing Satellites (AEOS) introduced new degrees of freedom in roll, pitch, and yaw. While this agility greatly enhances a satellite's ability to capture targets, it also exponentially increases the search space for potential schedules, making the optimization problem significantly more difficult \citep{Lemaitre2002}. Similarly, the recent proliferation of large satellite constellations has escalated the problem from a single-machine task to a far more complex multi-satellite coordination challenge, requiring new techniques to manage inter-satellite dynamics and data sharing \citep{augenstein2016optimal, Chen2019}. Each technological leap designed to improve data acquisition capability simultaneously imposes a greater computational burden on the scheduling systems that manage them.

This evolution in hardware has necessitated a parallel evolution in how the scheduling problem itself is conceptualized. Early models often focused narrowly on the primary mission of image acquisition \citep{Lemaitre2002,Gabrel2003}. However, it quickly became apparent that data downlink is a critical bottleneck, leading to the development of "integrated scheduling" models that consider both imaging and data transmission concurrently \citep{Wang2011,Hu2019,Zhang2021}. This study argues for a further evolution of this paradigm, from integrated mission planning to a more holistic form of resource lifecycle management. The operational reality of a satellite is a tightly coupled, cyclical process encompassing three key stages: image capture, data download, and on-board memory erasure. Effective scheduling cannot treat these as independent activities. An image capture consumes finite memory, which necessitates a timely download. Critically, the download itself does not free the on-board resource; an explicit delete command must be executed to reclaim that memory. This represents a complete \textit{consume $\rightarrow$ transfer $\rightarrow$ release} lifecycle for the memory resource, revealing hidden temporal and resource dependencies often overlooked in existing literature.

From an operational planning perspective, this setting is not only a combinatorial optimization problem but also a decision-support task in which heterogeneous mission requests, coupled resource states, and time-critical operational rules must be coordinated quickly and consistently. Effective planning therefore requires a scheduling framework that can combine explicit domain knowledge, optimization-based feasibility control, and fast heuristic search to produce implementable schedules under realistic satellite-control constraints.

The academic literature reflects a continuous effort to develop more powerful and realistic solutions. Early on, exact algorithms were employed to pursue provably optimal solutions. \citet{Gabrel2003} modeled a simplified version of the problem using graph theory, while more recent works have used sophisticated mathematical programming techniques like Mixed-Integer Linear Programming (MILP), Branch and Bound, and Branch and Price to handle greater complexity \citep{Chu2017,Chen2019,Hu2019,Peng2020,peng2025branch}. Despite their mathematical rigor, the primary drawback of these exact methods is their computational cost. The NP-hard nature of the scheduling problem renders them intractable for the large-scale instances and tight time constraints faced in daily satellite operations.

The scalability limitations of exact methods prompted a widespread shift toward heuristics and metaheuristics, which find high-quality, feasible solutions within practical computational limits. Seminal work by \citet{Lemaitre2002} investigated greedy algorithms and local search for agile satellites. Subsequently, a rich variety of metaheuristics has been applied, including Genetic and Evolutionary Algorithms \citep{Tangpattanakul2012,Tangpattanakul2013,Tangpattanakul2015,Zhang2021, wu2022data}, Ant Colony Optimization \citep{Wang2009, zhou2023multi}, Tabu Search \citep{Lin2005, sarkheyli2013using}, and Adaptive Large Neighborhood Search \citep{he2018improved, wu2024improved}. Recent studies have further explored advanced optimization variants, including differential-evolution ensemble learning for agile satellite scheduling \citep{chen2025deem}, multi-strategy NSGA-II for multi-objective scheduling of multiple agile satellites observing area targets \citep{mu2026multi}, and time-freedom-heuristic-guided ALNS for relay satellite scheduling with multi-type task requirements \citep{gu2026ensemble}. While these methods provide the necessary scalability, their lack of optimality guarantees create a fundamental trade-off. More recently, advances in machine learning have spurred interest in dynamic and autonomous scheduling using Deep Reinforcement Learning (DRL) and hybrid learning-search schemes, enabling satellites to react to new information in real time \citep{wei2021deep, herrmann2023reinforcement, qi2025deep, ahmadi2025enhancing}.

To navigate the trade-off between optimality and tractability, decomposition has emerged as a dominant strategic paradigm. Rather than solving a monolithic problem, these methods partition it into smaller, more manageable subproblems that can be solved more efficiently \citep{Boerkoel2021}. For instance, \citet{Wu2022} proposed a divide-and-conquer framework that first uses a metaheuristic to allocate tasks to orbits and then employs an exact solver to find the optimal schedule for each orbit. This hybridizes the global search power of heuristics with the local optimality of exact methods.

Despite this progress, a synthetic review of the state-of-the-art identifies several persistent research gaps:
\begin{enumerate}
    \item \textbf{Incomplete Mission Lifecycle Modeling:} The prevailing "integrated scheduling" paradigm is largely confined to image capture and data download \citep{Bianchessi2008,Hu2019}. It systematically overlooks the crucial final step of the on-board memory resource cycle: the explicit deletion of data. This omission leads to an inaccurate representation of memory availability and misses the complex temporal dependencies linking downloads to deletions and, consequently, to the capacity for new captures.
    \item \textbf{Insufficient Constraint Fidelity:} While individual studies have addressed various constraints like memory \citep{Wang2011} or power \citep{Lin2005}, few models holistically integrate the full, interdependent suite of operational rules. A comprehensive model must simultaneously manage on-board memory capacity, power consumption within orbital duty cycles, non-overlapping execution rules, and complex time-window dependencies dictated by ground station contact availability.
\end{enumerate}

This study addresses these gaps by developing a lifecycle-aware scheduling framework for remote-sensing satellite mission planning. We first formulate a mixed-integer model that captures the coupled capture-download-delete process, ground-station coordination, memory evolution, and orbital duty-cycle constraints. We then design a three-stage hybrid optimization architecture in which mission feasibility is determined first, playback sequencing is generated through a knapsack-to-shortest-path mechanism, and precise execution times are finalized through linear programming. This staged structure combines model-based rigor with efficient search, making the framework suitable for operational decision support. Computational experiments and sensitivity analyses show that the proposed framework produces high-quality feasible schedules within practical runtimes and can support planning decisions related to memory provisioning, contact availability, and download capability. The method has also been implemented in satellite-control operations, demonstrating practical applicability.

The core challenge is to maximize satellite-utilization efficiency with limited resources. We aim to maximize mission accommodation by considering the complete capture-download-delete lifecycle while adhering to a comprehensive suite of real-world operational constraints. The constructed mathematical model represents a complex MILP problem. To solve this, we propose a novel three-stage decomposition heuristic that separates the problem into distinct stages of mission selection, sequencing, and precise timing determination.
\begin{enumerate}
    \item \textbf{Stage 1 (Selection):} Determines mission-execution decisions while temporarily disregarding temporal factors.
    \item \textbf{Stage 2 (Sequencing):} Sequences selected missions by formulating them as multiple knapsack problems, which are then converted into shortest-path problems \citep{Dijkstra1959} for efficient optimization.
    \item \textbf{Stage 3 (Timing):} Solves a linear programming problem to determine precise mission-execution times.
\end{enumerate}
This approach generates high-quality, feasible schedules within reasonable computational times, balancing mission priorities and ensuring timely mission completion. Through empirical studies, we demonstrate that the proposed heuristic generates practical satellite schedules that satisfy operational requirements. The algorithm has been successfully implemented and deployed in satellite-control operations, validating its effectiveness and practical viability in real-world scenarios.


% MATHEMATICAL FORMULATION
\section{Mathematical Formulation}


This study addresses the comprehensive lifecycle management of satellite mission operations encompassing image acquisition (RSI record), data transmission (RSI playback), and memory resource release (RSI delete). The operational complexity stems from tightly coupled interdependencies: each RSI record generates multiple associated missions that must execute in coordinated sequences, ground station communications require precise attitude maneuvers (GAT for target acquisition, Gohome for power management), and memory resources follow a consume-transfer-release cycle creating temporal dependencies across mission groups.

The optimization challenge involves maximizing mission execution under stringent resource constraints: finite on-board memory capacity, limited filename allocation, restricted ground station contact windows (X-band for data transmission, S-band for status monitoring), and orbital power limitations within duty cycles. These constraints create a complex mixed-integer programming problem where mission selection decisions interact with precise temporal scheduling requirements, necessitating the comprehensive mathematical framework developed herein.


\subsection{Modeling Assumptions and Operational Framework}

This study adopts the following modeling assumptions to establish the mathematical framework scope:

\textbf{Modeling Assumptions:}
\begin{itemize}
    \item \textbf{Mission timing flexibility:} RSI playback and RSI delete missions have flexible execution windows defined by earliest and latest possible start times, with exact timing determined by the optimization model. All other missions have predetermined fixed execution schedules.
    \item \textbf{Ground station contact:} Satellite-ground station contact windows are precomputed based on orbital mechanics and remain fixed during the scheduling horizon, with stations classified by operational frequency (X-band or S-band).
    \item \textbf{Resource capacity constraints:} On-board memory capacity and filename allocation limits represent fixed system parameters that constrain mission execution feasibility.
\end{itemize}

\textbf{Operational Framework:}
The scheduling system operates under the following operational constraints that reflect real-world satellite mission requirements:
\begin{itemize}
    \item \textbf{Resource lifecycle management:} Memory resources follow a consume-transfer-release cycle where RSI record missions consume memory, RSI playback enables data transfer, and RSI delete operations explicitly release memory for reuse.
    \item \textbf{Power and orbital constraints:} Each orbital duty cycle (penumbra-to-penumbra interval) enforces maximum cumulative execution time limits for energy-intensive operations.
    \item \textbf{Temporal separation requirements:} Mission types with potential resource conflicts must maintain non-overlapping execution periods to prevent operational interference.
    \item \textbf{Mission interdependencies:} Certain mission combinations require coordinated execution, such as GAT and Gohome maneuvers accompanying ground station communications, and mandatory RSI delete execution following RSI record operations.
\end{itemize}

\begin{longtable}{|p{0.95\textwidth}|}
\caption{Operational constraints for satellite mission scheduling organized by constraint category}
\label{tab:constraints} \\
\hline
\textbf{TEMPORAL COORDINATION CONSTRAINTS} \\
\hline
\endfirsthead
\caption[]{Operational constraints for satellite mission scheduling (continued)} \\
\hline
\endhead
\hline
\endfoot
1. \textbf{Non-overlapping execution:} Mission execution periods within the same category cannot overlap (monitoring, attitude maneuvers, GAT maneuvers, Gohome maneuvers, RSI record, RSI playback, RSI delete). \\
2. \textbf{Cross-category conflicts:} Attitude, GAT, and Gohome maneuvers cannot overlap with each other or with RSI record/playback operations. \\
3. \textbf{Contact window alignment:} Monitoring missions must execute within S-band ground station contact periods; RSI playback missions must execute within X-band ground station contact periods. \\
4. \textbf{Sequential dependencies:} RSI delete operations can only commence after completion of all corresponding RSI record and playback missions in the same group. \\
5. \textbf{Timing buffers:} RSI delete execution must maintain minimum 6-second delay after the latest ground station contact end time for associated playback operations. Figure \ref{fig:rsitiming} illustrates this timing requirement where RSI delete can only start at least 6 seconds after the end of contact ID 80214 ($T \geq 6$). \\
\hline
\textbf{RESOURCE CAPACITY CONSTRAINTS} \\
\hline
6. \textbf{Memory management:} RSI record execution is constrained by available on-board memory capacity, with memory consumption determined by photography mode and release dependent on corresponding RSI delete execution. \\
7. \textbf{Filename allocation:} Total number of active image filenames cannot exceed system limits, with filename resources occupied from RSI record execution until corresponding RSI delete completion. \\
8. \textbf{Orbit-based power limits:} Cumulative execution time per orbital duty cycle cannot exceed 8 minutes for RSI record operations and 10 minutes for RSI playback operations. \\
\hline
\textbf{OPERATIONAL COORDINATION CONSTRAINTS} \\
\hline
9. \textbf{Ground station coordination:} RSI playback execution requires corresponding GAT and Gohome maneuvers for the selected ground station contact, with playback timing constrained between GAT completion and Gohome initiation. \\
10. \textbf{Mission group integrity:} If any RSI record mission in a group executes, the corresponding RSI delete mission must also execute to ensure memory resource release. \\
11. \textbf{Conditional maneuver execution:} GAT and Gohome maneuvers with specified contact IDs must execute simultaneously when the associated ground station is selected for RSI playback operations. \\
12. \textbf{Mission precedence:} RSI playback missions can only execute if the corresponding RSI record mission in the same group has been executed. \\
\hline
\end{longtable}
Based on the above framework, Table 1 organizes the operational constraints by constraint category:

\begin{figure}[h!]
    \centering
    \includegraphics[width=\columnwidth]{fig1.pdf}
    \caption{Illustrative example of time requirements of RSI delete operation}
    \label{fig:rsitiming}
\end{figure}

\subsection{Sets, Parameters, and Decision Variables}
This section introduces the sets, parameters, and decision variables used in the mathematical model.

\begin{longtable}{p{0.15\textwidth} p{0.8\textwidth}}
\caption{Sets, Parameters, and Decision Variables} \label{tab:notation} \\
\hline
\multicolumn{2}{l}{\textbf{Sets}} \\
\hline
\endfirsthead
\caption[]{Sets, Parameters, and Decision Variables (continued)} \\
\hline
\endhead
\hline
\endfoot
$A$ & Mission set with subsets $A', A'' \subseteq A$ for constraint formulation. \\
$N$ & Monitoring missions. $N \subseteq A$. \\
$T$ & Attitude maneuver missions. $T \subseteq A$. \\
$B$ & X-band ground-station contact groups indexed by contact ID. \\
$G$ & Gohome maneuver missions. $G \subseteq A$. \\
$C$ & Gohome maneuvers with specified contact IDs. $C = C^b = \{C^1, C^2, \dots\} \subseteq G \subseteq A$, $b \in B$. \\
$L$ & Gohome maneuvers without specified contact IDs. $L \subseteq G \subseteq A$. \\
$U$ & GAT maneuver missions. $U = U^b = \{U^1, U^2, \dots\} \subseteq A$, $b \in B$. \\
$H$ & Mission groups containing one RSI record, multiple RSI playbacks, and one RSI delete. \\
$R$ & RSI record missions. $R = R^h = \{R^1, R^2, \dots\} \subseteq A$, $h \in H$. \\
$P$ & RSI playback missions. $P = P_{i}^h = \{P_{1}^1, P_{2}^1, \dots, P_{1}^2, P_{2}^2, \dots\} \subseteq A$, $h \in H$, $i$ is activity index. \\
$D$ & RSI delete missions. $D = D^h = \{D^1, D^2, \dots\} \subseteq A$, $h \in H$. \\
$J$ & Ground stations categorized by frequency spectrum. \\
$X$ & X-band ground stations. $X = X^b = \{X^1, X^2, \dots\} \subseteq J$, $b \in B$. \\
$S$ & S-band ground stations. $S \subseteq J$. \\
$E$ & Eclipse periods. \\
$W$ & Memory and filename time windows based on RSI record execution timing, spanning from one record start to the next. \\
$Y$ & Duty-cycle time windows from penumbra entry to penumbra entry, constraining orbit-based execution limits. \\
\hline
\multicolumn{2}{l}{\textbf{Parameters}} \\
\hline
$M$ & Large constant for constraint formulation. \\
$t_i$ & Start time of mission $i \in A$ (earliest possible time for RSI playback and delete; exact time for others). \\
$d_i$ & Duration of mission $i \in A$. \\
$e_i$ & End time of mission $i \in A$ (latest possible time for RSI playback and delete; exact time for others). \\
$p_i$ & Priority weight of mission $i \in A$ for execution sequencing. \\
$r_i, f_i$ & Start and end times of ground station contact for $i \in X \cup S$. \\
$\eta_i, \xi_i$ & Start and end times of memory/filename time window $i \in W$. \\
$\mu$ & Initial available memory capacity. \\
$\omega_i$ & Memory consumed by RSI record $i \in R$. \\
$\delta_i$ & Memory released by RSI delete $i \in D$. \\
$\epsilon$ & Initial available filename count. \\
$\theta_i$ & Filenames occupied by RSI record $i \in R$. \\
$\lambda_i$ & Filenames released by RSI delete $i \in D$. \\
$\pi_i, \tau_i$ & Start and end times of duty-cycle window $i \in Y$. \\
\hline
\multicolumn{2}{l}{\textbf{Decision Variables}} \\
\hline
$v_i \in \{0,1\}$ & Mission execution indicator for $i \in A$. \\
$a_{ij} \in \{0,1\}$ & Temporal precedence indicator: $a_{ij} = 1$ if mission $i$ precedes mission $j$. \\
$b_{ij} \in \{0,1\}$ & Temporal precedence indicator: $b_{ij} = 1$ if mission $j$ precedes mission $i$. \\
$\sigma_{ij} \geq 0$ & Time offset from earliest possible start for mission $i$ using ground station $j \in X$ (RSI playback) or general offset for RSI delete. \\
$o_i \geq 0$ & Execution delay for flexible missions $i \in P \cup D$, set to large value for unexecuted missions. \\
$\phi_{ij} \in \{0,1\}$ & Ground station selection: $\phi_{ij} = 1$ if mission $i$ uses ground station $j \in X \cup S$. \\
$x_{ij} \in \{0,1\}$ & Time window assignment: $x_{ij} = 1$ if mission $i$ starts within window $j \in W \cup Y$. \\
\hline
\end{longtable}

\subsection{Mathematical model for satellite mission scheduling}

\begin{equation}
f_1(x) = \sum_{i \in A} p_i \times v_i \label{eq:obj1}
\end{equation}

\begin{equation}
f_2(x) = -\sum_{i \in P} p_i \times o_i - \sum_{i \in D} o_i \label{eq:obj2}
\end{equation}

\begin{equation}
\text{Maximize} \quad \psi f_1(x)+f_2(x) \label{eq:obj_combined}
\end{equation}

The mathematical model developed in this study has two main objectives: Equation~(\ref{eq:obj1}) prioritizes high-importance missions using priority weights $p_i$, accounting for overlapping tasks and limited resources. Equation~(\ref{eq:obj2})) minimizes $o_i$, the delay between each mission’s actual execution time and its earliest possible start. To further prioritize urgent image downloads, the first term in Equation~(\ref{eq:obj2}), $\sum_{i \in P} p_i \times o_i$, weights each RSI playback mission by its priority. Since the goal is to maximize the overall objective, this term appears with a negative sign. The optimization model integrates two complementary objectives through a weighted combination (Equation~(\ref{eq:obj_combined})). Parameter $\psi$ provides the necessary scaling to balance these objectives with different units and variable types. The constraints used in the model are described below:

\subsubsection{Constraints for nonoverlapping missions $A'$ and $A$}
\begin{equation}
t_i+\alpha+d_i-(t_j+\beta) \le M(1-a_{ij}), \quad \forall i \in A' \subseteq A, \forall j \in A'' \subseteq A \label{eq:nonoverlap1}
\end{equation}

\begin{equation}
t_j+\beta+d_j-(t_i+\alpha) \le M(1-b_{ij}), \quad \forall i \in A' \subseteq A, \forall j \in A'' \subseteq A \label{eq:nonoverlap2}
\end{equation}

\begin{equation}
v_i+v_j \le a_{ij}+b_{ij}+1, \quad \forall i \in A' \subseteq A, \forall j \in A'' \subseteq A \label{eq:nonoverlap3}
\end{equation}

The above three constraints ensure that missions cannot overlap during execution. Here, we use $\alpha$ and $\beta$ to represent the time difference between the actual and earliest possible start time for missions $i$ and $j,$ respectively. The values of $\alpha$ and $\beta$ will vary according to the types of relative missions $A'$ and $A''$ that need to be regulated, and the corresponding variable values will be inserted as presented in Table~\ref{tab:alpha_beta}.

\begin{table}[h!]
\caption{Values of $\alpha$ and $\beta$ for various nonoverlapping missions}
\label{tab:alpha_beta}
\centering
\small
\begin{tabular}{p{0.45\textwidth}ccp{0.25\textwidth}}
\toprule
\textbf{Type of Mission} & \textbf{$\alpha$} & \textbf{$\beta$} & \textbf{$i, j$} \\
\midrule
Mission start and end times are uncertain, but not affected by X-band ground-station contact times (i.e., RSI delete). & $\sigma_i$ & $\sigma_j$ & $\forall i, j \in D, i<j$ \\
Missions with definite execution start and end times. & 0 & 0 & $\forall i, j \in A' \subseteq A \setminus \{D, P\}, i<j$ \\
Missions with uncertain start and end times that are affected by X-band ground-station contact times (i.e., RSI playback). & $\sum_{k \in X} \sigma_{ik}$ & $\sum_{k \in X} \sigma_{jk}$ & $\forall i, j \in P, i<j$ \\
Mission $i$'s execution start and end times are definite, while the corresponding mission $j$ (RSI delete) has uncertain execution time but is not affected by X-band ground-station contact time. & 0 & $\sigma_j$ & $\forall i \in A' \subseteq A \setminus \{D,P\}, \forall j \in D$ \\
Both missions have definite execution times. & 0 & 0 & $\forall i \in A', \forall j \in A'', A', A'' \subseteq A \setminus \{D, P\}, A' \neq A''$ \\
Mission $i$'s execution start and end times are definite, while the corresponding mission $j$ (RSI playback) has an uncertain execution time and is affected by X-band ground-station contact time. & 0 & $\sum_{k \in X} \sigma_{jk}$ & $\forall i \in A' \subseteq A \setminus \{D,P\}, \forall j \in P$ \\
Both missions have uncertain start and end times, and only mission $j$'s execution time is affected by X-band ground-station contact time. & $\sigma_i$ & $\sum_{k \in X} \sigma_{jk}$ & $\forall i \in D, \forall j \in P$ \\
\bottomrule
\end{tabular}
\end{table}

\noindent The temporal offset parameters $\alpha$ and $\beta$ represent timing adjustments for missions with flexible start times: $\sigma_i$ for RSI delete missions, $\sum_{k \in X} \sigma_{ik}$ for RSI playback missions (with exactly one nonzero term), and 0 for fixed-time missions. The temporal precedence variables $a_{ij}$ and $b_{ij}$ encode relative mission timing: when missions overlap, exactly one variable activates to enforce mutual exclusion through constraint (\ref{eq:nonoverlap3}), limiting execution to at most one conflicting mission ($v_i + v_j \leq 1$). For non-overlapping missions, the optimization preference for maximum mission execution drives $a_{ij} = 1$ when mission $i$ precedes $j$, allowing concurrent execution ($v_i + v_j \leq 2$).

\subsubsection{Type $A'$ mission and ground station $J'$ contact-time constraint}
\begin{equation}
t_i+\alpha \ge r_j \times \phi_{ij}, \quad \forall i \in A' \subseteq A, \forall j \in J' \subseteq J \label{eq:contact1}
\end{equation}

\begin{equation}
t_i+\alpha+d_i \le f_j \times \phi_{ij} + M(1-\phi_{ij}), \quad \forall i \in A' \subseteq A, \forall j \in J' \subseteq J \label{eq:contact2}
\end{equation}

\begin{equation}
\sum_{j \in J' \subseteq J} \phi_{ij}-v_i = 0, \quad \forall i \in A' \subseteq A \label{eq:contact3}
\end{equation}

\begin{equation}
\sigma_{ij} \le M\phi_{ij}, \quad \forall i \in P, \forall j \in X \label{eq:contact4}
\end{equation}

These constraints enforce mission execution within ground station contact windows. Constraints (\ref{eq:contact1})-(\ref{eq:contact2}) ensure mission timing falls within selected contact intervals $[r_j, f_j]$ through binary variable $\phi_{ij}$. Constraint (\ref{eq:contact3}) requires exactly one contact selection per executed mission, while constraint (\ref{eq:contact4}) maintains variable consistency by nullifying timing variables for unselected ground stations.

\subsubsection{Memory and number of filenames constraint}
\begin{equation}
t_i+o_i \ge \eta_j \times x_{ij}, \quad \forall i \in D, \forall j \in W \label{eq:memory1}
\end{equation}

\begin{equation}
t_i+o_i+d_i \le \xi_j \times x_{ij} + M(1-x_{ij}), \quad \forall i \in D, \forall j \in W \label{eq:memory2}
\end{equation}

\begin{equation}
v_i \times \omega_i \le \mu - \sum_{k \in \{R|k<i\}} v_k \times \omega_k + \sum_{l \in D} \sum_{j \in \{W|j<i\}} \delta_l \times x_{lj}, \quad \forall i \in R \label{eq:memory3}
\end{equation}

\begin{equation}
v_i \times \theta_i \le \varepsilon - \sum_{k \in \{R|k<i\}} v_k \times \theta_k + \sum_{l \in D} \sum_{j \in \{W|j<i\}} \lambda_l \times x_{lj}, \quad \forall i \in R \label{eq:memory4}
\end{equation}

These constraints manage dynamic memory and filename allocation through time-windowed resource tracking. Constraints (\ref{eq:memory1})-(\ref{eq:memory2}) assign RSI delete operations to temporal windows based on their execution timing. Constraints (\ref{eq:memory3})-(\ref{eq:memory4}) enforce capacity limits by tracking cumulative resource consumption from executed RSI records minus releases from completed deletions, ensuring resource availability throughout the mission lifecycle.

\subsubsection{Duty-cycle time constraint}
\begin{equation}
t_i \times v_i + \alpha \ge \pi_j \times x_{ij}, \quad \forall i \in A' \subseteq A, \forall j \in Y \label{eq:duty1}
\end{equation}

\begin{equation}
t_i \times v_i + \alpha \le \tau_j \times x_{ij} + M(1-x_{ij}), \quad \forall i \in A' \subseteq A, \forall j \in Y \label{eq:duty2}
\end{equation}

\begin{equation}
v_i - \sum_{j \in Y} x_{ij} = 0, \quad \forall i \in A' \subseteq A \label{eq:duty3}
\end{equation}

\begin{equation}
\sum_{i \in A' \subseteq A} d_i \times x_{ij} \le \text{time}, \quad \forall j \in Y \label{eq:duty4}
\end{equation}

These constraints enforce orbital power limitations by restricting cumulative execution time within duty cycles. Constraints (\ref{eq:duty1})-(\ref{eq:duty2}) assign missions to time windows $[\pi_j, \tau_j]$ based on execution timing, while constraint (\ref{eq:duty3}) ensures each executed mission belongs to exactly one duty cycle. Constraint (\ref{eq:duty4}) enforces cycle-specific execution time limits to maintain power consumption within operational bounds.

\subsubsection{RSI-delete time constraint}
\begin{equation}
t_i+\sigma_i \ge \sum_{k \in X} f_k \times \phi_{jk} + \text{time}, \quad \forall i \in D^h, \forall j \in P^h, \forall h \in H \label{eq:delete1}
\end{equation}

\begin{equation}
\sigma_i \le M v_i, \quad \forall i \in D \label{eq:delete2}
\end{equation}
\begin{equation}
o_i = \sigma_i + M(1-v_i), \quad \forall i \in D \label{eq:delete3}
\end{equation}

These constraints manage RSI delete timing requirements. Constraint (\ref{eq:delete1}) enforces minimum 6-second delay between delete execution and the latest ground station contact end time for associated playback operations. Constraints (\ref{eq:delete2})-(\ref{eq:delete3}) maintain timing variable consistency, setting $\sigma_i = 0$ for unexecuted missions while assigning penalty values for objective function evaluation.

\subsubsection{RSI-playback time constraint}
\begin{equation}
o_i = \sum_{j \in X} \sigma_{ij} + M(1-v_i), \quad \forall i \in P \label{eq:playback1}
\end{equation}

\begin{equation}
t_i+\sigma_{ik}+M(1-\phi_{ik}) \ge e_j, \quad \\ \forall i \in P, \forall j \in U^b, \forall k \in X^b, \forall b \in B \label{eq:playback2}
\end{equation}

\begin{equation}
t_i+\sigma_{ik}+d_i-M(1-\phi_{ik}) \le t_j, \quad \\ \forall i \in P, \forall j \in C^b, \forall k \in X^b, \forall b \in B \label{eq:playback3}
\end{equation}

These constraints coordinate RSI playback execution with ground station selection and associated maneuvers. Constraint (\ref{eq:playback1}) maintains timing consistency for playback operations. Constraints (\ref{eq:playback2})-(\ref{eq:playback3}) ensure RSI playback execution occurs between corresponding GAT maneuver completion and Gohome maneuver initiation when a ground station is selected.

\subsubsection{Mission correlation constraints}
\begin{equation}
v_i-\phi_{jk} \ge 0, \quad \forall i \in U^b, \forall j \in P, \forall k \in X^b, \forall b \in B \label{eq:corr1}
\end{equation}

\begin{equation}
v_i-v_j = 0, \quad \\ \text{for } (\forall i \in U^b, \forall j \in C^b, \forall b \in B) \text{ or } (\forall i \in R^h, \forall j \in D^h, \forall h \in H) \label{eq:corr2}
\end{equation}

\begin{equation}
v_i-v_j \ge 0, \quad \forall i \in R^h, \forall j \in P^h, \forall h \in H \label{eq:corr3}
\end{equation}
These constraints enforce operational dependencies between related missions. Constraint (\ref{eq:corr1}) requires GAT maneuver execution when the corresponding ground station is selected for RSI playback. Constraint (\ref{eq:corr2}) enforces simultaneous execution of GAT and Gohome maneuvers for the same contact, and mandatory RSI delete execution when RSI record is performed. Constraint (\ref{eq:corr3}) ensures RSI playback can only execute if the corresponding RSI record in the same group has been executed.

\section{Three-Stage Decomposition Heuristic}
While the mathematical model accurately captures the satellite scheduling
  problem, its NP-hard nature renders direct solution computationally intractable
  for real-time operational requirements, with solution times growing exponentially
   with problem size. To ensure scalability and solution quality, we develop a tailored three-stage decomposition heuristic. Section~\ref{pc} analyzes problem characteristics, and Section~\ref{ds} outlines the decomposition logic designed to efficiently generate near-optimal solutions through subproblem solving.

\subsection{Problem characteristics}\label{pc}
In this satellite mission scheduling problem, many mission types are time-deterministic missions, which only require judging whether the mission can be executed based on mission constraints. This selective problem requires choosing executable missions without violating feasibility to maximize the objective value. However, missions such as RSI playback and delete require the model to determine the exact start time of mission execution. Further, RSI playback is affected by multiple factors, such as whether the GAT and Gohome maneuvers are executed, contact time between the satellite and ground station, and duration of contact, these influence its exact execution time. 

Furthermore, there can be various possible combinations of RSI playback and ground stations. Let us assume $k$ RSI playback type missions that must be scheduled, where each RSI playback in the initial input data provides only the earliest and latest possible start times for execution. Within this interval, the satellite will pass over multiple ground stations. If all RSI playbacks can choose among $s$ ground stations, and each RSI playback must specify one ground station for downloading, there will be $s^k$ possible combinations. As the number of RSI playbacks increases, the total number of combinations increases exponentially, making the solution complex. Finally, the start time of execution also affects the performance of the objective value, making the RSI playback mission the bottleneck in the overall solution. To solve this problem, we design a decomposition strategy, breaking down the solution process into multiple stages.

\subsection{Decomposition strategy}\label{ds}
The proposed decomposition strategy is shown in Figure~\ref{fig:heuristic_process}. The first stage selects as many potential missions as possible without violating constraints. The second stage determines the optimal execution order for bottleneck missions, specifically RSI playbacks, to complete them as early as possible. The third stage solves the original model with the given order and fixed variables to obtain exact execution times. If no feasible solution is found due to the imposed order, additional variables are introduced to relax these constraints and ensure a feasible and effective result.

\begin{figure}[h!]
    \centering
    \includegraphics[width=0.8\textwidth]{fig2.pdf}
    \caption{Three-Stage decomposed heuristic solution process}
    \label{fig:heuristic_process}
\end{figure}

\subsubsection{Selection of missions to execute}
The primary goal of this satellite mission scheduling problem is to ensure that all missions in the scheduling list are executed, as far as possible, and then start the RSI playback and deletion with uncertain execution start and end times at the earliest. The first step of the proposed solution scheme ignores the time objective of starting early and only maximizes the sum of executable mission weights (i.e., to solve the model with objective function~(\ref{eq:m1_obj}) (complete solution model shown as $\textbf{M1}$)). Preliminary experiments revealed that solving this model can result in convergence to the near-optimal solution in a short time, and the selected missions to be executed can serve as input for the second step.

\noindent\textbf{Model M1:}
\begin{equation}
\max f'(x) = \sum_{i \in A} p_i \times v_i \label{eq:m1_obj}
\end{equation}
subject to Constraints \eqref{eq:nonoverlap1}-\eqref{eq:corr3}

\subsubsection{Execution order determination for RSI playback missions}
Due to varying contact durations and earliest execution times, each RSI playback mission may be assigned to multiple ground stations, resulting in high computational complexity in the original model. Preliminary tests show that the execution order of RSI playbacks significantly affects scheduling quality. Different sequences can lead to different completion times. To improve performance, we designed a sequencing mechanism that prioritizes missions with higher priority and earlier start times, and schedules them as early as possible to better utilize ground station availability. Starting with the first ground station, missions are filtered by priority and start time. If all selected missions can be scheduled, their order is determined. Otherwise, unscheduled missions are combined with the remaining ones and reassessed for the next station. This process continues until all missions are considered. 

The problem for each ground station is converted into a knapsack problem, where each time segment downloadable by a ground station is viewed as a knapsack. The total available download time corresponds to the total weight capacity of the knapsack, while the duration of each RSI playback represents the weight of items to be placed in the knapsack. Additionally, because the goal is to fill the knapsack as much as possible, the duration of each RSI playback also serves as the value of the item. The objective is to maximize the total value that can be placed in the knapsack, i.e., to minimize the remaining space. Using Figure~\ref{fig:knapsack} as an example, we assume that a ground station has a downloadable time segment of 50 s, with three RSI playback missions to consider (durations of 20, 10, and 30 s each). We can draw a network diagram accordingly. If RSI Playback 1 is chosen to be executed at this ground station, we traverse from the starting point $\textbf{S}$ to node $1^2$, with a corresponding edge cost (benefit) of -20. Otherwise, we traverse from $\textbf{S}$ to node $1^0$, with a corresponding edge cost (benefit) of 0. RSI Playback 2 and RSI Playback 3 use the same logic to construct the network. Additionally, a network endpoint $\textbf{T}$ is added. Under this design, any path from starting point $\textbf{S}$ to endpoint $\textbf{T}$ represents a placement of knapsack items (downloading RSI playback missions at this ground station) that does not violate the knapsack capacity constraint (the 50 s downloadable time segment limit of the ground station). Because the edge costs correspond to negative time benefits, by determining the shortest path from starting point $\textbf{S}$ to endpoint $\textbf{T}$, the corresponding solution will be the mission execution method with maximum benefit that does not violate the downloadable time-duration regulations of the ground station.

\begin{figure}[h!]
    \centering
    \includegraphics[width=0.8\textwidth]{fig3.pdf}
    \caption{Conversion of the knapsack problem into a shortest path problem}
    \label{fig:knapsack}
\end{figure}

Note that converting the problem into a knapsack problem helps select a scheduling approach that maximizes the utilization of ground-station availability. Specifically, the following are the heuristic steps designed to achieve the abovementioned goal:

\begin{description}
    \item[Step 0:Initialization of parameters :] Including the ground-station set $CON$, which contains all downloadable ground stations operational during the scheduling period; the pending RSI playback set $PBK$, which stores all executable RSI playbacks from the first-stage scheduling results; and the ordered RSI playback set $ORD$, which stores the final RSI playbacks arranged in execution order. This list is initially empty.
    \item[Step 1:] The ground stations in $CON$ are sorted in chronological order based on their start times.
    \item[Step 2:] The earliest ground station is selected from the sorted $CON$ and the RSI playbacks in $PBK$ are examined. The RSI playbacks that should be considered together for sequencing are chosen. The method for selecting which RSI playbacks should be grouped together for sequencing will be described in the next paragraph. If no RSI playback can be executed at the selected ground station, remove that ground station from $CON$ and proceed to the next suitable one. If there is no next ground station, skip to \textbf{Step 5}.
    
    When examining RSI playbacks, both their priority and the earliest possible execution time must be considered. In this study, Equation~(\ref{eq:rank}) classifies missions into ranks based on these criteria. Missions assigned the same rank are scheduled together in a coordinated sequence. In this equation, $p$ represents the priority of the RSI playback, while $ss$ denotes the index of the ground station capable of downloading the mission. For example, if three ground stations can download a mission, the corresponding $ss$ values would be 1, 2, and 3, respectively. A larger $ss$ value indicates that the download is deferred to a later ground station; for example, $ss=1$ implies that the RSI playback selects the first ground station in its candidate list.

    \begin{equation}
    rank = 3^{p} \times ss^3. \label{eq:rank}
    \end{equation}
    
    \item[Step 3:] Based on the selected RSI playback missions, the designed shortest path problem is solved using the Dijkstra algorithm \citep{Dijkstra1959}.
    \item[Step 4:] The RSI playback missions selected from the shortest path problem in \textbf{Step 3} for the current ground station are sorted and inserted into the list $ORD$ in the following order: from the highest to lowest priority, from the earliest to latest possible start time, and from the shortest to longest execution duration. These missions are then removed from $PBK$. If the ground station has no remaining available download time after this step, it is removed from $CON$.
    \item[Step 5:] The termination condition satisfaction (i.e., either $CON$ or $PBK$ is empty) is checked. If the condition is not yet satisfied, return to \textbf{Step 2}; otherwise, terminate the process.
\end{description}

\subsubsection{Determination of precise mission execution time}
To determine the final execution time of each mission, we additionally design constraint~(\ref{eq:order_const}):

\begin{equation}
t_i + \sum_{k \in X} \sigma_{ik} \le t_j + \sum_{k \in X} \sigma_{jk} + M(2-v_i-v_j) \quad \forall i, j \in P, i<j. \label{eq:order_const}
\end{equation}

In this equation, suppose that, after the aforementioned steps, RSI playback $j$ is scheduled after RSI playback $i$. Then, the start time of RSI playback $i$ (i.e., $t_i + \sum_{k \in X} \sigma_{ik}$) must be earlier than the start time of RSI playback $j$ (i.e., $t_j + \sum_{k \in X} \sigma_{jk} + M(2-v_i-v_j)$). This ensures that precise execution start times can be assigned under the given order. In other words, model $\textbf{M2}$ is solved to determine the final mission execution times.

\noindent\textbf{Model M2:}
\begin{align}
\max \quad & \text{Eq.}~(\ref{eq:obj_combined}) \nonumber
\end{align}
subject to constraints (\ref{eq:nonoverlap1})–(\ref{eq:corr3}) and (\ref{eq:order_const})

\subsubsection{Final adjustment}
In the proposed third-stage strategy, constraint~(\ref{eq:order_const}) is not considered during the first stage when mission executability is determined. The second stage sets mission order, but during the third stage, some missions initially marked as executable may become infeasible due to constraints like memory limits. Since maximizing the number of executed missions is the primary goal, the model restores feasibility by relaxing execution order constraints for RSI playback when necessary. For this purpose, a new binary variable $\rho_{ij} \in \{0, 1\}$ is introduced and incorporated into the original constraint~(\ref{eq:order_const}), resulting in a new constraint~(\ref{eq:order_relax}):

\begin{equation}
t_i + \sum_{k \in X} \sigma_{ik} \le t_j + \sum_{k \in X} \sigma_{jk} + M(2-v_i-v_j) + M\rho_{ij} \quad \forall i, j \in P, i<j \label{eq:order_relax}
\end{equation}

In this equation, if RSI playback $i$ is originally required to precede RSI playback $j$ according to the sequencing result from the second stage, but the order is violated in the final schedule determined in the third stage, then $\rho_{ij}=1$; otherwise, $\rho_{ij}=0$. Furthermore, to minimize the deviation between the final schedule and optimized order obtained in the second stage, this relaxation variable is included in the objective function, as presented in Equation~(\ref{eq:m3_obj}) (with the complete model denoted as $\textbf{M3}$). This penalizes violations of the stage-two sequencing rules and guides the solution to remain as close as possible to the initial optimized order. A sensitivity analysis was conducted to test and calibrate the appropriate weight value $\chi$; Section \ref{senpsichi} presents the details.

\noindent\textbf{Model M3:}
\begin{equation}
    \max f''(x) = \psi \sum_{i \in A} p_i \times v_i + (-\sum_{i \in P} p_i \times o_i - \sum_{i \in D} o_i) + (-\chi \sum_{i, j \in P} \rho_{ij}) \label{eq:m3_obj}
\end{equation}
subject to constraints (\ref{eq:nonoverlap1})–(\ref{eq:corr3}) and (\ref{eq:order_relax})

Based on the above design, the original mixed-integer programming problem is effectively simplified into a series of near-linear programming problems, significantly reducing the computational complexity.

\section{Numerical Experiments}
To evaluate the effectiveness of the proposed scheduling framework, we implemented it in Java and tested it on a workstation with an Intel i7-11700 CPU @2.50GHz processor and 56 GB memory. The experiments are designed to assess three aspects of the framework: solution quality, computational efficiency under operational time limits, and decision-support value through sensitivity analysis.

\subsection{Determination of parameters $\psi$ and $\chi$} \label{senpsichi}
Parameters $\psi$ and $\chi$ are introduced to balance different objectives within the optimization models. Specifically, $\psi$ is used to scale and integrate two objective functions with different units and variable types, binary in the first and continuous in the second, in the original mathematical formulation. Meanwhile, $\chi$ penalizes violations of the sequencing results obtained from a prior optimization stage, encouraging alignment between intermediate and final solutions in the third stage of the proposed solution strategy. Because the appropriate values of $\psi$ and $\chi$ directly affect the performance and trade-off behavior of the model, sensitivity analysis is necessary to test and calibrate these parameters in numerical experiments. These tests were validated using real data provided by the collaborating organization. Based on our numerical tests, the parameter setting of $\psi=10^6$ and $\chi=1$ yields the most balanced and practically applicable performance. Therefore, we adopt these two values in the subsequent experiments.

\subsection{Computational evaluation}
To assess the computational performance of the proposed framework, we conducted experiments on eight mission sets of increasing size, summarized in Table~\ref{tab:instances}. Each result is averaged over 10 runs to reduce the effect of algorithmic variability. Because satellite mission planning is subject to strict operational response limits, CPLEX was used as a time-limited benchmark with a total solving budget of 6 minutes, divided into two 3-minute stages corresponding to mission selection and RSI playback scheduling. Since CPLEX may not reach optimality within this limit, we report its best incumbent solutions for comparison. The objective of this experiment is not only to compare solution quality, but also to evaluate whether the proposed framework can generate high-quality feasible schedules within runtimes suitable for operational decision support.

\begin{table}[h!]
\caption{Problem instances for validation}
\label{tab:instances}
\centering
\begin{tabular}{lcccc}
\toprule
 & \# Record & \# Playback & \# Delete & \# Missions \\
\midrule
Case 1 & 1 & 5 & 1 & 24 \\
Case 2 & 2 & 10 & 2 & 33 \\
Case 3 & 3 & 15 & 3 & 43 \\
Case 4 & 4 & 20 & 4 & 52 \\
Case 5 & 5 & 25 & 5 & 60 \\
Case 6 & 8 & 40 & 8 & 88 \\
Case 7 & 12 & 57 & 12 & 122 \\
Case 8 & 17 & 81 & 17 & 167 \\
\bottomrule
\end{tabular}
\end{table}

\noindent Table~\ref{tab:validation} shows that the proposed framework remains competitive with CPLEX on small instances and becomes increasingly advantageous as problem size grows. For Cases 1--3, both methods produce fully feasible schedules with similar quality, although the proposed framework already offers substantially shorter runtimes in the larger of these small cases. From Case 4 onward, the benefits of the staged optimization structure become clearer. In medium and large instances, the proposed method schedules more missions, achieves smaller completion-time delays, and does so within significantly lower computational times than time-limited CPLEX. Case 8, which corresponds to the real-world scale used in system development, particularly highlights the framework's suitability for operational planning.

\begin{table}[h!]
\caption{Validation results}
\label{tab:validation}
\centering
\begin{tabular}{lcccccc}
\toprule
 & \multicolumn{3}{c}{CPLEX} & \multicolumn{3}{c}{Proposed Method} \\
\cmidrule(lr){2-4} \cmidrule(lr){5-7}
 & NM* & T ** & CPU*** & NM* & T ** & CPU*** \\
\midrule
Case 1 & 0 & 3,475.90 & 1.32 & 0 & 3,475.90 & 1.75 \\
Case 2 & 0 & 4,061.30 & 1.56 & 0 & 4,060.50 & 2.21 \\
Case 3 & 0 & 7,442.10 & 111.19 & 0 & 7,440.10 & 2.45 \\
Case 4 & 0 & 5,894.70 & 332.31 & 0 & 5,894.30 & 3.15 \\
Case 5 & 0.50 & 7,064.40 & 362.06 & 0.50 & 7,064.40 & 39.99 \\
Case 6 & 3.10 & 13,406.90 & 362.69 & 0.50 & 12,881.80 & 47.74 \\
Case 7 & 5.40 & 34,289.70 & 363.47 & 0.80 & 32,754.40 & 47.94 \\
Case 8 & 14.80 & 46,754.20 & 364.88 & 0.40 & 39,543.70 & 108.41 \\
\bottomrule
\end{tabular}
\\
*Number of unscheduled missions ;  
**Time difference (in seconds) between the completion time of the final playback mission and earliest available download time; 
*** Solution time in seconds
\end{table}

\subsection{Sensitivity analysis}
Beyond solution quality, the proposed framework is intended to support planning decisions under changing system conditions. To examine this decision-support capability, we conducted sensitivity analyses on satellite memory capacity, ground-station service time, and image download speed. These factors are operationally important because they affect both scheduling feasibility and infrastructure requirements. By varying them systematically, we evaluate how the framework responds to changes in resource availability and communication capability. The results, averaged over five randomly generated cases, are summarized in Table~\ref{tab:sensitivity}.

\subsubsection{Sensitivity analysis: Memory size}
As shown in Table~\ref{tab:sensitivity}, increasing memory capacity significantly improves scheduling performance. With more available memory, the number of unscheduled missions decreases, especially for image download missions such as RSI playback. When memory is limited, captured images are often deleted before being fully downloaded, making these missions more prone to disruption. Once the memory reaches 3,400 Mbit or above, all missions can be successfully scheduled. This sensitivity analysis highlights the algorithm’s ability to evaluate the relationship between mission volume and memory requirements. For instance, when scheduling approximately 160 missions, including 16 sets of image acquisition tasks, a memory capacity of at least 3,400 Mbit is recommended to ensure both scheduling feasibility and computational efficiency.

\subsubsection{Sensitivity analysis: Available time}
As a satellite orbits the Earth, it establishes brief communication links with multiple ground stations to download images. Due to maintenance or competing satellite requests, a station may offer only partial service during its contact window, requiring the satellite to use additional stations. These contact durations, typically 3 to 8 minutes, play a critical role in scheduling decisions. To assess the impact of ground-station contact duration, we varied the available service time of one randomly selected station and observed changes in overall scheduling performance. As shown in Table~\ref{tab:sensitivity}, when the station provides no service (available time = 0 s), the satellite requires five additional stations to meet download demands, denoted as “1 + 5” (one unavailable plus five others). As contact duration increases, the number of required stations gradually decreases. Once the duration exceeds 3 minutes (available time $\ge$ 180 s), the satellite can complete all downloads using only five stations.

\begin{table}[h!]
\caption{Sensitivity analysis}
\label{tab:sensitivity}
\centering
\begin{tabular}{llcccc}
\toprule
 & & NM* & NS** & CPU*** & Opt. Gap \\
\midrule
 & 1,000 & 13.8 & 5 & 279.57 & 2.76\% / 0.01\% \\
 & 1,600 & 9.8 & 4 & 240.44 & 4.47\% / 0\% \\
 & 2,200 & 4.2 & 5.8 & 333.87 & 2.44\% / 0.05\% \\
\textbf{Memory Size (Mbit)} & 2,800 & 1.8 & 5 & 188.80 & 1.04\% / 0.03\% \\
 & 3,400 & 0 & 5 & 75.81 & 0\% / 0\% \\
 & 4,000 & 0 & 5 & 54.24 & 0\% / 0\% \\
\midrule
 & 0 & 0 & 1+5 & 64.50 & 0\% / 0\% \\
 & 60 & 0 & 6 & 81.73 & 0\% / 0\% \\
 & 120 & 0 & 5.6 & 65.33 & 0\% / 0\% \\
\textbf{Available Time (s)} & 180 & 0 & 5 & 51.28 & 0\% / 0\% \\
 & 240 & 0 & 5 & 34.50 & 0\% / 0\% \\
 & 300 & 0 & 5 & 57.03 & 0\% / 0\% \\
\midrule
 & 1 & 0 & 5 & 41.65 & 0\% / 0\% \\
 & 0.9 & 0 & 5 & 26.07 & 0\% / 0\% \\
\textbf{Download Speed} & 0.7 & 0 & 4 & 39.95 & 0\% / 0\% \\
\textbf{(measured as a multiple} & 0.5 & 0 & 4 & 52.04 & 0\% / 0\% \\
 \textbf{of the playback duration)}& 0.3 & 0 & 4 & 37.02 & 0\% / 0\% \\
 & 0.1 & 0 & 4 & 28.20 & 0\% / 0\% \\
\bottomrule
\end{tabular}
\\
*NM: Number of unscheduled missions; **NS: Total number of ground stations used; ***CPU: Solution time in seconds
\end{table}

\subsubsection{Sensitivity analysis: Download speed}
The image download rate influences both download time and ground-station service costs. While current systems use the X band, emerging technologies like optical terminals may increase speeds and reduce reliance on ground stations. This analysis adjusts download times to fractions of the original (1, 0.9, 0.7, 0.5, 0.3, 0.1) to assess the impact on ground-station demand. As shown in Table~\ref{tab:sensitivity}, at the original speed, five ground stations are required. When download time is reduced to 0.7 or less, only four stations are needed. This indicates that once download speed surpasses a certain threshold, ground-station demand drops notably. Fewer ground stations lower communication costs, increase efficiency, and allow more imagery to be downloaded within the same window, enabling faster and more cost-effective mission execution.

\section{Conclusion}
This study addresses the computationally challenging problem of remote-sensing satellite mission scheduling by introducing the first comprehensive framework to model the complete capture-download-delete mission lifecycle, thereby accurately representing on-board memory resource dynamics previously overlooked in existing integrated scheduling paradigms. We develop a novel three-stage decomposition heuristic that systematically balances mission feasibility, resource optimization, and timing requirements through mission selection, knapsack-to-shortest-path sequencing, and precise temporal assignment via linear programming. Our approach demonstrates competitive and scalable performance across problem instances of varying complexity (24-167 missions), achieving near-optimal solutions with substantially reduced mission infeasibility relative to time-limited commercial solver benchmarks while maintaining computational efficiency suitable for real-time operational deployment. Comprehensive sensitivity analyses reveal critical threshold behaviors for system design: minimum memory capacity requirements (3,400 Mbit for large-scale operations), optimal ground station utilization patterns, and technology upgrade decision points that enable informed infrastructure investment strategies. The methodology has been successfully implemented and validated in operational satellite control systems, demonstrating both theoretical rigor and practical applicability for industrial-scale mission scheduling. Future research directions include multi-satellite constellation coordination with inter-satellite data sharing, integrated ground station network optimization, and robust scheduling frameworks incorporating stochastic operational uncertainties to further enhance the framework's resilience and expand its applicability to next-generation satellite mission management systems.

\section*{Acknowledgments}
The authors acknowledge the National Science and Technology Council, Taiwan, ROC, for providing partial funding support under contract numbers NSTC 112-2410-H-007-060-MY3 and NSTC 114-2410-H-007-057-MY3. The contents of the article remain the sole responsibility of the authors.

\bibliographystyle{elsarticle-harv}

\bibliography{cas-refs.bib}

\end{document}

